\documentclass[10pt,a4paper]{article}

%========================
% 中文与字体
%========================
\usepackage[UTF8]{ctex}

%========================
% 页面设置
%========================
\usepackage[
    left=1.6cm,
    right=1.6cm,
    top=1.3cm,
    bottom=1.3cm
]{geometry}

%========================
% 常用包
%========================
\usepackage{enumitem}
\usepackage{titlesec}
\usepackage{hyperref}
\usepackage{xcolor}


%========================
% 基础设置
%========================

\setlength{\parindent}{0pt}

\hypersetup{
    colorlinks=true,
    urlcolor=black
}


% bullet间距
\setlist[itemize]{
    leftmargin=1.5em,
    itemsep=0pt,
    topsep=1pt,
    parsep=0pt
}


% section样式
\titleformat{\section}
{\large\bfseries}
{}
{0em}
{}
[\titlerule]


\titlespacing{\section}
{0pt}
{6pt}
{4pt}



\begin{document}


%========================
% 个人信息
%========================

\begin{center}

{\fontsize{18pt}{20pt}\selectfont
\textbf{孙一鸣}
}

\vspace{6pt}

(+86) 157-1288-7586
\quad | \quad
\href{mailto:1120233185@bit.edu.cn}{1120233185@bit.edu.cn}
\quad | \quad
\href{https://github.com/melquias-s}{GitHub: melquias-s}

\end{center}



%========================
\section{教育背景}
%========================


\textbf{北京理工大学机电学院, 机械电子工程（全英文）}
\hfill
2023.09 - 至今


\begin{itemize}

\item 
GPA: 3.6/4.0；专业排名：7/17；CET-6:557（2024）

\item 
核心课程：
机电控制技术（99）、
卡尔曼滤波与惯性导航基础（95）、
机器人学（93）、
计算机技术与编程（91）、
机械原理（90）

\end{itemize}




%========================
\section{科研与项目经历}
%========================


% -------- 项目1 --------

\textbf{仿生扑翼机器人跳跃起飞机构设计}

\begin{itemize}

\item
面向仿生扑翼机器人起飞需求，设计轻量化弹性储能跳跃腿机构，实现机械能储存与瞬时释放驱动的快速起跳。

\item
基于仿生学分析与信鸽运动特征研究，利用 SolidWorks 完成机构设计，并结合 3D 打印与实物测试进行迭代优化。

\item
设计由弹性骨架、驱动绳和线性/非线性弹簧组成的储能机构，并设计电机驱动的能量耦合/解耦锁定机制。

\item
完成约 50 g 轻量化样机，实现稳定 2.5 m 高度跳跃。

\end{itemize}



% -------- 项目2 --------

\textbf{基于视觉感知的麦克纳姆轮自主机器人系统}

\begin{itemize}

\item
设计并实现自主机器人视觉感知系统，基于 Python/OpenCV 实现 HSV 颜色分割、轮廓检测、ArUco 定位及目标坐标计算。

\item
基于 K-Means 聚类提升多目标识别与定位稳定性，并设计目标丢失补偿机制，提高复杂环境下视觉鲁棒性。

\item
搭建视觉模块与 Raspberry Pi 控制端通信流程，实现机器人坐标、目标位置及路径点实时传输。

\item
负责机器人系统方案设计、结构优化与调试，实现视觉、控制和机械模块协同运行。

\end{itemize}



% -------- 项目3 --------

\textbf{仿生飞鱼机器人结构重建与三维建模}

\begin{itemize}

\item
基于 CT 扫描数据，利用 3D Slicer 完成生物结构分割与三维重建，建立仿生飞鱼机器人结构数字模型。

\item
提取关键形态特征并优化模型结构，为仿生机器人设计提供数据支持。

\item
完成适用于 3D 打印制造的模型转换，为机器人实体原型制造提供基础。

\end{itemize}

% -------- 项目4 --------

\textbf{水空两栖跨介质机器人嵌入式控制系统开发}

\begin{itemize}

\item
基于 STM32 开发机器人嵌入式控制模块，实现基于 PWM 调制的舵机、电机控制、UART 通信及蓝牙无线通信。

\item
设计基于帧协议的数据传输机制，实现控制指令解析与可靠通信。

\item
开发遥控舵机测试 Demo，实现舵机归中、角度循环摆动及旋钮实时控制。

\item
研究基础飞控控制逻辑，包括 V 型尾翼运动控制与执行机构混控策略。

\end{itemize}




%========================
\section{技术能力}
%========================

\begin{itemize}

\item
编程语言：Python、C/C++、MATLAB

\item
机器人感知：OpenCV、计算机视觉、目标识别与定位

\item
嵌入式系统：STM32、Raspberry Pi、嵌入式通信、执行机构控制

\item
三维建模与制造：SolidWorks、3D Slicer、MeshLab、3D Printing

\end{itemize}




%========================
\section{荣誉奖项}
%========================

\begin{itemize}[leftmargin=1.5em]

\item 北京理工大学五四优秀共青团员、本科生二等奖学金
\hfill 2026

\item 北京理工大学五四优秀共青团员、优秀实践团员、本科生进步奖学金
\hfill 2025

\item 北京理工大学优秀班干部、本科生三等奖学金
\hfill 2024

\item 美国大学生数学建模竞赛（MCM/ICM）Honorable Mention
\hfill 2025

\end{itemize}



%========================
\section{学生工作}
%========================

\begin{itemize}

\item
校级学生组织智能无人系统队队长

\item
机械电子工程（全英文）班班长

\item
精工书院第三党支部副书记

\end{itemize}



\end{document}